\documentclass[sigconf,review]{acmart}

\AtBeginDocument{%
  \providecommand\BibTeX{{%
    Bib\TeX}}
}

% TODO: Replace with final ACM metadata after acceptance, if required.
\settopmatter{printacmref=false}
\renewcommand\footnotetextcopyrightpermission[1]{}
\pagestyle{plain}

\title{Smart Coder: Dynamic Autonomy for Coding Agents Under Explicit Local Governance}

% TODO: Replace author metadata before submission.
\author{Tanja L. Shukla}
\email{TODO}
\affiliation{%
  \institution{TODO}
  \city{TODO}
  \state{TODO}
  \country{USA}
}

\renewcommand{\shortauthors}{Shukla}

\begin{document}

\begin{abstract}
Coding agents force an uncomfortable tradeoff: either developers approve nearly every action, which slows work, or they grant broad autonomy and accept the risk of subtle bugs, unsafe changes, and scope drift. We present Smart Coder, a governance-first framework for adaptive autonomy in coding agents. Smart Coder treats the language model as untrusted: the model may request reads, declare intent, initiate check-ins, and propose file updates, but it cannot directly read files, write files, or execute commands. A local CLI mediates every side effect using hard constraints, workflow phase, verification outcomes, and interaction traces. The system combines two independent oversight sources: policy-initiated check-ins from the CLI and model-initiated check-ins when the model surfaces uncertainty, assumptions, or architectural tradeoffs. Every interaction is logged and reused to adapt both future policy decisions and future prompt context. The current prototype implements a heuristic baseline with milestone-aware oversight, read/write-split constraints, spec-aware planning, and exportable study traces. We position Smart Coder as both a working research artifact and a framework for studying how adaptive autonomy can better balance developer productivity and safety in agent-assisted software engineering.
\end{abstract}

\keywords{coding agents, adaptive autonomy, human-AI collaboration, software engineering, AI oversight, agentic systems}

\maketitle

\section{Introduction}

Coding agents now generate substantial code and execute multi-step workflows, but the harder problem is no longer only generation quality; it is oversight calibration. Prior work on agent oversight, appropriate reliance, and human-centered evaluation all points to the same tension: too much approval friction collapses productivity, while too much unchecked autonomy increases the risk of subtle bugs, security mistakes, and over-trust~\cite{grunde2026overseeing,bansal2021appropriate,wang2025humansmissing}. Recent field observations of professional developers reinforce that this is not a niche concern: effective use of coding agents remains highly controlled in practice, with developers deliberately steering planning, verification, and risky implementation choices rather than delegating end-to-end autonomy~\cite{huang2025dontvibe}.

Developers currently navigate this problem with a mix of markdown rule files, manual review, narrower prompts, and personal habits for when an agent should pause. These workarounds help, but they remain fragmented. Static rule files can state some explicit preferences, while repeated corrections and approval patterns reveal others only through use. Existing coding-agent workflows offer limited support for turning both kinds of signals into persistent, inspectable oversight behavior without collapsing everything into either prompt text or model-internal memory.

We present \emph{Smart Coder}, a governance-first framework for adaptive autonomy in coding agents. A formative survey with 21 experienced coding-agent users reinforced the need for milestone- and risk-based check-ins, visible reasons for changing autonomy, and less interruption on routine work. Smart Coder responds by keeping authority in a local CLI while adapting from explicit user-authored rules and observed interaction history. Its central contributions are a governed runtime in which the model remains structurally untrusted, a dual-initiator check-in design spanning both the model and the policy layer, and a trace-driven adaptation loop that changes future prompting and future oversight across sessions.

\section{System Overview}

\begin{figure}[t]
  \centering
  \fbox{\parbox{0.95\columnwidth}{\textbf{Placeholder for architecture figure.} Replace this box with the Smart Coder architecture diagram: developer $\leftrightarrow$ CLI governance layer $\leftrightarrow$ untrusted model, with SQLite traces, policy engine, prompt builder, and verification loop.}}
  \caption{Smart Coder architecture. The model proposes actions; the local CLI remains the sole enforcement boundary for file access, writes, and verification.}
  \label{fig:architecture}
\end{figure}

Smart Coder is implemented as a local Python CLI that sits between the developer and a Bedrock-hosted coding model. The model operates through a strict structured protocol: it may request file reads, declare intent, issue check-ins, and propose file updates, but it cannot directly read or modify the repository. The CLI validates every request against repository-local state and decides whether to continue, prompt, or block.

This design can also be understood through a capability-security lens~\\cite{lampson1971protection}. Smart Coder's leases are not generic trust toggles; they are temporary, path-scoped capabilities for a specific access type, with expiry and local revocation.

\subsection{Governance and adaptive autonomy}

Every read and write request passes through an approval cascade. First, Smart Coder checks hard constraints, which encode deterministic read/write policies such as \texttt{always\_deny}, \texttt{always\_check\_in}, and \texttt{always\_allow}. Second, it checks active leases, which are temporary trust grants derived from prior approvals. Third, it evaluates a heuristic policy score built from interaction history and risk signals. The score is based on factors such as prior approvals and denials, review speed, edit distance, blast radius, change type, diff size, security sensitivity, and verification-failure history. The user does not directly manipulate this scoring logic; instead, they interact through qualitative autonomy modes such as \emph{strict}, \emph{balanced}, \emph{milestone}, and \emph{autonomous}.

This policy layer is deliberately hybrid rather than purely heuristic. Some checkpoints are deterministic and milestone-driven regardless of score: plan approval, phase transitions, constrained files, deviations from an approved plan, and verification hotspots can all trigger review. This hybrid design matches the interaction style preferred by many of our survey respondents: do not interrupt for routine work, but do stop at meaningful decision points.

\subsection{Dual-initiator check-ins}

Smart Coder supports two independent sources of oversight. The first is policy-initiated oversight from the CLI. The second is model-initiated oversight, where the model is explicitly instructed to surface uncertainty, assumptions, and design tradeoffs through structured \texttt{CheckInMessage} outputs. This matters because not all valuable pauses can be expressed as static rules. A model may recognize that preserving an API response envelope conflicts with a cleaner abstraction, or that a spec appears to require an interface change. When that happens, Smart Coder allows the model to ask for guidance rather than silently choose.

Each check-in is logged with an explicit initiator field. This makes the system auditable and creates a natural research question for future study: whether CLI-side or model-side check-ins are better calibrated for different users or tasks.

\subsection{Trace-driven prompt and policy loop}

Smart Coder's core mechanism is an external trace loop. Every interaction produces structured traces that include the stage, file path, change type, risk features, user decision, review timing, free-text feedback, verification result, and model metadata such as self-reported confidence and check-in assumptions. These traces are persisted in SQLite and reused in two places.

First, traces affect future governance decisions. Repeated denials, high edit distance, verification failures, or low-confidence behavior can tighten autonomy in a file or pattern. Repeated successful approvals can loosen it. Second, traces are summarized into future system prompts. The model sees qualitative trust summaries, relevance-ranked historical corrections, behavioral guidelines, autonomy preferences, and workflow phase. It does not see raw internal thresholds or exact trust scores.

This externalized trace-to-policy and trace-to-prompt loop is a central distinction between Smart Coder and both static configuration files and model-internal personalization methods such as PAHF-style memory updates~\\cite{liang2026pahf}. In the current prototype, prompt-side retrieval already goes beyond simple recency: historical corrections and accepted guidelines are ranked against the current task and spec context before being injected. The adaptation remains inspectable, local, and separable from the underlying model.

\subsection{Spec-aware planning and verification}

Smart Coder optionally accepts a task specification via \texttt{--spec}. The spec is digested into prompt context, and the model is asked to map its plan to requirements and declare potential deviations. The CLI then treats missing requirement coverage or explicit deviations as reasons to stop at the plan gate. This produces a lightweight spec-driven workflow without requiring a full external planning system.

After approval, updates are applied atomically and passed through a configured verification command. Verification outcomes are logged and fed back into future autonomy decisions. This matters for safety: repeated verification failures should reduce autonomy even if a developer has previously approved similar edits quickly.

\section{Demo Scenario and Evidence}

The demo uses a small external Python task API repository with route handlers, a service layer, structured application errors, tests, and a short task specification. The specification acts as a task contract: it preserves response envelopes and handler signatures while leaving implementation strategy open. Before coding begins, the developer authors two freeform rules with \texttt{sc rules add}: one that compiles into a CLI-enforced constraint and one that compiles into prompt-level behavioral guidance. This setup makes the separation between deterministic governance, softer preferences, and task-specific specification visible before the main coding task starts.

The live workflow then proceeds in two short sessions. In session 1, Smart Coder receives the task and the spec, produces a spec-aware plan, and surfaces an API design tradeoff through a model-initiated check-in rather than guessing silently. The developer chooses an option and gives one explicit instruction about how subsequent low-risk work should proceed. In session 2, Smart Coder performs a related follow-up task on the same repository, now starting from loaded interaction history and reducing unnecessary interruption while preserving review at the API boundary. The demo closes with one observability view, showing that the run was fully traced and exported. This sequence highlights the system's intended contribution: adaptive oversight, not static prompting or unconstrained autonomy.

\begin{table}[t]
  \caption{Evidence to include in the final submission from one or more representative rehearsal runs.}
  \label{tab:evidence}
  \centering
  \begin{tabular}{p{0.44\columnwidth}p{0.42\columnwidth}}
    \toprule
    Evidence item & Current source \\
    \midrule
    Hard-rule compilation example & \textbf{TODO:} \texttt{sc rules add} output \\
    Guidance-rule compilation example & \textbf{TODO:} \texttt{sc rules add} output \\
    Representative task success & \textbf{TODO:} exported session bundle \\
    Session 1 model vs policy check-ins & \textbf{TODO:} \texttt{observe checkin-stats} \\
    Session 2 fewer interruptions than Session 1 & \textbf{TODO:} session trace comparison \\
    Session 2 starts from loaded history & \textbf{TODO:} run bootstrap output \\
    Verification outcomes & session bundle + trace CSV \\
    \bottomrule
  \end{tabular}
\end{table}

At minimum, the final submission should include one clean exported session bundle, one trace CSV, and a short video demonstrating the full two-session flow. Those artifacts support the three claims this paper needs: the system runs end to end on a realistic coding task, both model-side and policy-side oversight occur, and observed developer behavior changes later system behavior.

\section{Conclusion and Limitations}

Smart Coder's main contribution is not a claim of fully solved personalization; it is a governed adaptive-autonomy framework that makes oversight behavior explicit, persistent, and observable. By separating task contracts from softer preferences and by logging both model-initiated and policy-initiated check-ins, the system creates a concrete platform for studying how coding-agent autonomy should adapt over time rather than treating that behavior as an opaque model property.

The current prototype remains a heuristic baseline. Its approval policy uses manually chosen priors rather than learned calibration, its spec support is lightweight rather than fully structured, and it does not yet implement richer interrupt taxonomy, post-hoc correction, rewind, or asynchronous delegation. These are meaningful limitations, but they are also appropriate for a demo paper whose goal is to present a working artifact and an instrumented workflow rather than a finalized learning system.

The immediate next step is a controlled lab study using the packaged CLI, the demo repository, and exported traces from repeated sessions. That study will let us contribute a clearer taxonomy of user preference types, test whether different oversight styles emerge across users and tasks, and evaluate whether learned policies or fine-tuned models can outperform the current heuristic runtime. Beyond coding, the same framework may generalize to other agentic workflows where user preferences mix hard governance rules with softer behavioral conditioning.


\bibliographystyle{ACM-Reference-Format}
\bibliography{refs}

\end{document}
